\section{Method}

\subsection{Task Definition}

We study \emph{verdict prediction}: given a Python function, a test input, and an optional execution trace, predict whether the test passes or fails.
The model produces a single token (\texttt{pass} or \texttt{fail}), making this a binary classification task.
Unlike code generation \citep{chen2025tracepile} or output prediction \citep{gu2024cruxeval}, verdict prediction isolates the model's ability to detect faults from its ability to simulate execution.
The discrete binary outcome makes the effect of each trace component directly measurable, providing cleaner ablation contrasts than generation tasks where output quality varies continuously.

\subsection{Dataset}
\label{sec:dataset}

We sampled 10 candidate solutions per problem from Qwen-2.5-Coder-7B-Instruct (temperature 0.8, top\_p 0.95) for 591 problems from HumanEval \citep{chen2021evaluating} (164) and MBPP-sanitized \citep{austin2021program} (427), executed each against canonical tests with \texttt{sys.settrace} to obtain line-level traces and pass/fail verdicts, yielding 5,910 pairs (67\% pass, 33\% fail).
The dataset is split 80/20 by problem ID into 4,720 training and 1,190 test samples (majority baseline: 71.51\%).
The \texttt{exception\_only} format collapses passing solutions into identical rows (39.45\% train, 48.99\% test duplicates); these duplicates arise from multiple candidate solutions per problem, with the train--test split enforced at the problem level (472 train / 119 test problems, zero crossover; Appendix~\ref{sec:dedup_stats}).

\subsection{Trace Representations}

We define seven conditions that vary the trace information available (Figure~\ref{fig:trace_overview}):
\begin{itemize}
\item \texttt{no\_trace}: Source code and test input only.
\item \texttt{full\_trace}: Complete line-by-line trace with variable states.
\item \texttt{VOTA}: Exception events + assertion-point comparisons + loop compression (Appendix~\ref{sec:vota_algorithm}).
\item \texttt{exception\_only}: Only exception/error events.
\item \texttt{ae\_only}: Only assertion-point actual/expected pairs.
\item \texttt{loop\_only}: Only loop compression summaries.
\item \texttt{VOTA\_no\_ae}: VOTA without assertion-point information.
\end{itemize}

\begin{figure*}[t]
\centering
\includegraphics[width=\textwidth]{fig1_trace_overview}
\caption{Overview of the verdict prediction pipeline. Source code and a test input are executed to produce a full raw trace. Four conditions vary the trace information: \texttt{no\_trace} (75.13\%), \texttt{full\_trace} (81.76\%), VOTA (99.41\%), and \texttt{exception\_only} (99.83\%). Accuracies for Qwen-2.5-7B (seed=42, $N{=}1{,}190$). Appendix~\ref{sec:vota_examples} provides input/output examples.}
\label{fig:trace_overview}
\end{figure*}

\subsection{Models and Training}

We evaluate four models spanning three architecture families: \textbf{Qwen-2.5-Coder-7B/1.5B-Instruct} \citep{qwen2025coder}, \textbf{CodeLlama-13B-Instruct} \citep{roziere2024code}, and \textbf{DeepSeek-Coder-6.7B-Instruct} \citep{deepseek2024coder}, covering a 1.5B--13B parameter range (8.7$\times$).
All conditions use LoRA \citep{hu2022lora} (rank=16, $\alpha$=32, dropout=0.05) on all linear layers, 3 epochs, effective batch size 32, LR $2 \times 10^{-4}$ with linear warmup/decay, max sequence length 2{,}048 tokens.
Sequences exceeding this limit are truncated while preserving the verdict label.
The primary seed is 42; seeds 123 and 456 test \texttt{exception\_only} robustness (Appendix~\ref{sec:training_details}).

\subsection{Loss Ablation}

To disentangle training signal from input signal quality, we introduce \texttt{full\_trace\_label\_only}: identical full-trace input but loss computed only on the verdict token (trace positions masked).
Improvement over standard \texttt{full\_trace} is attributable to \emph{gradient dilution}; the remaining gap to VOTA reflects \emph{information overload} at inference time.
